\documentclass[11pt]{article}
\usepackage{fontspec,amsmath,amssymb,amsthm,hyperref}
\usepackage[margin=1in]{geometry}
\setlength{\emergencystretch}{3em}
\title{Move the center of a finite projected RMS error}
\author{LangevinLean mathematical exposition}
\date{}
\begin{document}\maketitle
\noindent Source module: \texttt{ProximalLangevin/ProjectionCenter.lean}. Compile this proof document with XeLaTeX or LuaLaTeX. The accompanying HTML places the complete Lean declarations beside every proof.


For a probability measure \(\mu\) and bounded linear map \(T\), write \(R_\mu(a)=\|T(\cdot)-a\|_{L^2(\mu)}\), initially as an extended nonnegative number. Finiteness is checked before converting this norm to a real number.


\section{Minkowski and finiteness}


The identity \(Tx-a=(Tx-b)+(b-a)\) and Minkowski yield \(R_\mu(a)\le R_\mu(b)+\|b-a\|\). A constant has that \(L^2\) norm because \(\mu\) has mass one. If \(R_\mu(b)<\infty\), the right side is finite, hence the error about \(a\) is finite too.



\noindent Lean declarations: \texttt{projectedRMS\_center\_le}, \texttt{projectedRMS\_center\_finite}.

\section{Pass to ordinary square-root integrals}


Monotonicity of the extended-to-real conversion is valid once the upper bound is finite. This proves the real RMS inequality. Use the proved finite moments at both centers to identify the norms with \(\sqrt{\int\|Tx-a\|^2d\mu}\) and \(\sqrt{\int\|Tx-b\|^2d\mu}\). The resulting formula cannot hide an infinite moment through the totalized real integral.



\noindent Lean declarations: \texttt{projectedRMSReal\_center\_le}, \texttt{sqrt\_integral\_center\_le}.

\end{document}
